\subsubsection{Stray Light}
\label{ssec:stray}
%
The compact, 3-mirror design of the Simonyi Survey Telescope (\secref{ssec:simonyi}) is especially sensitive to stray and scattered light\footnote{Light that does not follow the nominal optical path within the telescope and instrument, originating from scattering off optical and non-optical surfaces or from ghost paths in refractive components such as lenses and filters}.
The most pronounced and extended artifacts appear in \gls{LSSTCam} images when the telescope's boresight is near bright objects such as the Moon, bright stars, planets, or unintended light sources on the sky.
Stray light contaminates all astronomical imaging to some level, but \gls{LSST}'s low-surface-brightness science is especially sensitive to it.

A significant fraction of these artifacts traced back to the Observatory dome, which lacked its \gls{LWS} to restrict the acceptance angle along the telescopes's elevation axis; many artifacts are expected to resolve once the \gls{LWS} is comissionined, whereas others will require additional baffling.
Some artifacts detected were expected based on an initial stray light analysis in FRED\registered \cite[][]{2009SPIE.7427E..08E}, which correctly predicted direct paths from the primary (M1) and tertiary (M3) mirrors that bypass the secondary mirror and hit the \gls{LSSTCam} optics and focal plane directly.
Other artifacts were new and unexpected, requiring visual inspection, ray tracing simulation, and on-sky tests to confirm their origins \citep[][]{2026arXiv2606.31945}.

An initial census of stray light artifacts observed in \gls{LSSTCam} images during initial commissioning (April 2025–May 2026) is reproduced in \tabref{tab:rates} from \cite[][]{2026arXiv2606.31945}.
About 70\% of features come from identified, confirmed sources, and 80\% already have a possible mitigation strategy.
Occurrence rates are based on visual classification of features visible by eye; for some, fractional area is a rough estimate due to variation across frames.

\begin{deluxetable*}{lcccc}
\tablecaption{Occurrence rate of stray and scattered light artifacts identifiable by visual inspection of $\sim$\,116,000 full focal plane mosaic images.
Rates are not expected to be complete nor pure; impacted area uses average feature size.
The most prevalent artifact, the \textit{scratched tape}, was effectively mitigated through installation of the mid-level baffle extension \citep[][]{2026arXiv2606.31939}.
\vspace{1em}
*\hspace{0.1cm}Fractional area of the artifact is variable; \dag \hspace{0.1cm} artifact origin to be confirmed; $\Upsilon$ preliminary estimate.  \label{tab:rates}}
\tablehead{\colhead{Artifact} & \colhead{Source} & \colhead{Mitigation} & \colhead{Occurrence Rate}  & \colhead{Fractional Area}}
\startdata
Scratched Tape & \gls{M1} direct path & Mid-Level Baffle Extension & 7.39\% & 2.2\% * \\
Smudge & \gls{M1} direct path & Mid-Level Baffle Extension & 0.19\% & 0.28\% \\
Muddy Shoe & \gls{M1} direct path & Mid-Level Baffle Extension & 0.15\% & 2\%* \\
Pillow & Unknown & \gls{LWS} installation & 0.26\% & 9\%* \\
Nibbled Disc & Unknown  & Mid-Level Baffle Extension & 0.10\% & 1.06 - 10\%* \\
L2 Bananas & L2 support pads & \gls{LWS} installation & 0.73\% & 0.7\% \\
Rivet Holes & Holes in Top-end Baffle & Sealing of holes & 0.45\% & -- \\
Sharp Arcs & \gls{L2} edges \dag & TBD & TBD & 1\% $\Upsilon$ \\
Submarine & \gls{L1}/L2 ghosts \dag & LWS installation (TBC) & 0.7\% $\Upsilon$ & 8.5\%* \\
\gls{L3} Horseshoe & \gls{L3} chamfer \& mount & Additional \gls{L3} edge baffle & 0.09\% & 8.8\% \\
\hline \\[-1em]
All artifacts & --- & --- & 10.00\% & --- \\
\multicolumn{2}{c}{Post Mid-Level Baffle Extension installation} & --- & 2.6\% & ---\\
\enddata
\end{deluxetable*}

The most frequent and impactful feature,  the \textit{scratched tape} feature \citep[][]{2026arXiv2606.31939}, originates from the M1 direct path that skips M2 and M3 reflections, and was visible in $\sim$\,7.4\% of image frames with an off-axis range of visibility between $\sim$\,19\fdeg8 and 22$^\circ$.
An annular extension of the telescope's mid-level baffle ring, installed between February and April 2026, elimintated this feature from \gls{LSSTCam} images \citep[][]{2026arXiv2606.31939}.
Other features such as the \textit{L3 Horseshoe}, arise from edge diffraction at \gls{LSSTCam}'s tertiary lens (L3) near the field-of-view edge ($\sim1\fdeg99$ and $2\fdeg23$).
Ray tracing simulations also revealed scattered light from the M2 optical baffle, whose residual reflectance in the longer wavelength bands ($i$, $z$, $y$) can reach up to 30\%.

Despite the preventative measures taken, $\sim$\,10\% of early image frames suffer from stray light contamination, dropping to $\sim$\,2.6\% following the mid-level baffle extension installation \citep[][]{2026arXiv2606.31939}.
The most severe issues were structured artifacts on the focal plane, largely stemming from the incomplete status of the \gls{LWS}.
 Rubin is implementing new physical baffling and re-painting to address the remaining artifacts \citep[][]{2026arXiv2606.31898, 2026arXiv2606.31793}.
